\section{Compact pairings and comparison coordinates}
\label{sec:comparison}

The construction in this section is topological.  A Riemannian metric is
used only to measure derivatives in the direction of the original
manifold.  In particular, no curvature hypothesis enters the following
statement. For a vector, $\norm{\cdot}_\infty$ is the maximum of the
absolute coordinate values. Bounds asserted on a set are pointwise;
uniform norms take the supremum over the domain.

\begin{theorem}[Comparison data from rational essentiality]
\label{thm:comparison-data}
Let $(M^n,g_M)$ be a closed connected oriented Riemannian manifold,
$n\ge2$, let $\Gamma=\pi_1(M)$, and suppose that its classifying map
$c_M:M\to B\Gamma$ satisfies
\[
 (c_M)_*[M]\ne0\quad\text{in }H_n(B\Gamma;\Q).
\]
There are constants $J>0$ and $R_0\ge1$ such that, for every real
$R\ge R_0$, there exist a closed oriented smooth manifold $A$, a
principal left $\Gamma$-cover $\widehat A\to A$, and an integer
$q=n+\dim A>n$ with the following properties.  The diagonal cover
\[
 Z=(\widetilde M\times\widehat A)/\Gamma\longrightarrow M\times A
\]
is oriented and has dimension $q$.  There are a smooth compactly
supported map $G:Z\to\R^q$ and a relatively compact open set
$\Omega\subset Z$ such that
\begin{equation}
\label{eq:comparison-degree}
 \deg(G,\Omega,0)\ne0,
 \qquad \norm{G}_\infty>3\quad\text{on }\partial\Omega,
\end{equation}
and
\begin{equation}
\label{eq:comparison-bounds}
 \max_{1\le i\le q}\sup_Z\abs{d_{\widetilde M}G_i}_{g_M}
 \le\frac JR,
 \qquad \log q\le JR.
\end{equation}
Write $\pi_Z:Z\to M\times A$ for the covering projection. Its
differential gives a global splitting
\[
 TZ=E_M\oplus E_A,\qquad
 E_M=(d\pi_Z)^{-1}(TM\oplus0),\qquad
 E_A=(d\pi_Z)^{-1}(0\oplus TA).
\]
The bundle $E_M$ is tangent to the $\widetilde M$ fibres of $Z\to A$,
and $E_A$ is the flat lift of the parameter directions. We write
$d_{\widetilde M}G_i=dG_i|_{E_M}$. For each fixed $R$, after choosing any
smooth metric on $A$, all derivatives of $G$ in the lifted product
geometry are bounded.  No bound uniform in $R$ is asserted for
parameter derivatives or for derivatives of order greater than one.
The covers and the parameter manifolds need not be connected.
\end{theorem}

We use rational coefficients for the chain constructions in this
section. A principal left $\Gamma$-cover is a covering on which $\Gamma$
acts freely by deck transformations and transitively on each fibre. Every
diagonal quotient below uses left actions in both factors.  Orientations
on covers are lifted from the base, so all deck transformations preserve
them.

The construction proceeds as follows. An auxiliary closed manifold $C$
provides a finite parameter space for a compactly supported group-cohomology
pairing. Distance profiles then supply $q$ scalar coordinates, with $q$
initially only a row count. Relative duality produces a complementary
rational class of degree $q-n$. Oriented realization of a nonzero multiple,
followed by doubling, gives the parameter manifold $A$. Thus the final cover has
exactly the dimension needed for a local degree, $n+\dim A=q$; the topology
of the realized $A$ does not alter the previously chosen row count.

\subsection{Probability maps and a finite parameter manifold}

Fix a finite symmetric generating set of $\Gamma$, and write
$d_\Gamma$ for its left-invariant word metric.  For $a\ge0$, let
$\mathcal K_a$ be the Rips complex realized as finite probability vectors
$t=(t_\gamma)_{\gamma\in\Gamma}$ whose nonempty support has diameter at
most $a$.  This complex is locally finite, finite dimensional, and
cocompact under left translation.  The action is proper, although it
need not be free when $\Gamma$ has torsion.

\begin{lemma}[A proper probability map]
\label{lem:probability-map}
There are $s\ge1$, $H<\infty$, a smooth equivariant probability map
$h:\widetilde M\to\mathcal K_s$, and an equivariant simplicial map
$h_0:\widetilde M\to\mathcal K_s$ such that
\begin{enumerate}
\item the coordinate functions of $h$ are smooth and locally supported
on finitely many labels, and $\abs{dh}_{\ell^1}\le H$;
\item $h$ and $h_0$ are joined by an equivariant homotopy that is proper
as a map $\widetilde M\times[0,1]\to\mathcal K_s$;
\item there is an invariant alternating homogeneous rational group
$n$-cocycle $\omega$ whose evaluation on the labelled fundamental
simplices of $M$ is nonzero.
\end{enumerate}
\end{lemma}

\begin{proof}
Choose a finite smooth triangulation of $M$ by the Cairns--Whitehead
theorem \cite{Cairns1934,Whitehead1940}; see also
\cite[Theorem~10.6]{Munkres1966}. Label the vertex lifts
equivariantly by elements of $\Gamma$, making one arbitrary choice for
each vertex downstairs.  Barycentric coordinates define the labelled
simplicial map $h_0$.  A smooth partition of unity subordinate to the
open vertex stars, lifted equivariantly, gives $h$; weights with the same
label are added.  At each point the labels that occur belong to the
vertices of one lifted simplex.  Cocompactness bounds their word
diameters by a fixed $s$, and it also bounds the $\ell^1$-operator norm
of $dh$ by a fixed $H$.  Linear interpolation of the weights stays in
$\mathcal K_s$.

For completeness, this entire homotopy is proper.  If $Q\subset
\widetilde M$ is compact with $\Gamma Q=\widetilde M$, its image during
the homotopy is compact.  Given a compact $B\subset\mathcal K_s$,
properness of the $\Gamma$-action allows only finitely many translates
of that image to meet $B$.  The inverse image of $B$ is therefore a
closed subset of finitely many translates of $Q\times[0,1]$.

Choose a class in $H^n(B\Gamma;\Q)$ pairing nontrivially with
$(c_M)_*[M]$, and represent it on the homogeneous bar resolution.  On
tuple chains, signed permutation averaging is an equivariant chain map
lifting the identity of the augmentation.  It is equivariantly chain
homotopic to the identity on this free resolution.  Antisymmetrizing
the representative thus gives a cocycle $\omega$ in the same class.
It vanishes on repeated labels and restricts to a simplicial cocycle on
every $\mathcal K_a$.

Give the vertices of $M$ a global order.  Sending an ordered lifted
simplex to its tuple of labels gives an augmentation-preserving
equivariant chain map into the bar resolution.  Such a map is unique
up to equivariant chain homotopy: inductively, the difference in each
degree is a cycle in the exact target resolution, and freeness of the
source chains permits equivariant fillings.  Comparison with the chain
map induced by a classifying map shows that evaluation of $\omega$ on
these labelled fundamental simplices is precisely its pairing with
$(c_M)_*[M]$, and hence is nonzero.
\end{proof}

The detecting cocycle is now fixed. To construct a finite smooth parameter
space carrying its compact pairing, we need a free auxiliary group action.
The ordinary Rips action may have torsion stabilizers; keeping the order
of every simplex removes them without requiring contractibility.

Fix henceforth $R\ge s+1$ and set
\begin{equation}
\label{eq:comparison-scales}
 L=100R.
\end{equation}
All auxiliary spaces constructed at this scale may depend on $L$.
Let $\mathcal O_L$ be the semisimplicial complex with one simplex for
each ordered tuple of distinct group elements of diameter at most $L$;
faces are obtained by deleting slots, and different orderings are not
identified.  Projection of a tuple to its probability simplex gives
$\pi:\mathcal O_L\to\mathcal K_L$.  On chains it has the equivariant
section
\begin{equation}
\label{eq:ordered-section}
 \sigma[v_0,\ldots,v_j]
 =\frac1{(j+1)!}\sum_{\varpi\in S_{j+1}}
    \sgn(\varpi)(v_{\varpi(0)},\ldots,v_{\varpi(j)}).
\end{equation}
Counting the positions of a deleted entry verifies that $\sigma$ is a
chain map, and $\pi_\#\sigma=\id$.

\begin{lemma}[A parameter thickening]
\label{lem:parameter-thickening}
Put $v_0=\dim\mathcal O_L$.  There are a closed oriented smooth
$d$-manifold $C$, with $d\le2v_0+2$, a principal left $\Gamma$-cover
$\widehat C\to C$, and equivariant maps
\[
 \mathcal O_L\xrightarrow{i}\widehat C
 \xrightarrow{r}\mathcal K_L,
 \qquad r i=\pi.
\]
Moreover $v_0\le\abs{B_\Gamma(1,L)}$.
\end{lemma}

\begin{proof}
A point in the realization of $\mathcal O_L$ belongs to the interior of
a unique minimal-dimensional ordered simplex.  A group element fixing
the point must preserve that ordered simplex, hence must fix its labels
slot by slot.  It is therefore the identity.  This also proves freeness
at points represented on faces; face identifications do not identify
distinct orderings of the minimal simplex.  Local finiteness and the
word-metric bound make the action proper and its quotient finite.
Thus $\mathcal O_L\to\mathcal O_L/\Gamma$ is a principal cover, including
when $\Gamma$ has torsion.  The quotient has a finite polyhedral
subdivision.  A simplex has at most $\abs{B_\Gamma(1,L)}$ vertices after
one of its vertices is translated to the identity, proving the
dimension bound.

Embed the finite polyhedron $\mathcal O_L/\Gamma$ in
$\R^{2v_0+1}$ and choose an open neighborhood retracting onto it.
Choose a compact smooth domain inside that neighborhood containing the
polyhedron in its interior, and double the domain.  Folding the double
to the domain followed by the neighborhood retraction gives a
retraction from a closed oriented manifold $C$ onto
$\mathcal O_L/\Gamma$.  Its dimension satisfies the stated bound.
Pull back the principal cover along this retraction.  The inclusion
of the polyhedron lifts to $i$, while the pullback projection followed
by $\pi$ gives $r$.  Their composition has the required value.
\end{proof}

Fix a smooth metric on $C$ and lift it to $\widehat C$.  A useful
consequence of continuity and cocompactness is the following uniform
common-label property: sufficiently small subsets of $\widehat C$ have
images under $r$ with a label common to every positive support.  Indeed,
at any point a positive coordinate remains positive on a neighborhood;
a finite cover downstairs gives a uniform neighborhood size upstairs.

\begin{lemma}[A controlled map to dual chains]
\label{lem:dual-carrier}
One can triangulate $C$ so finely that the lifted triangulation has
the following properties:
\begin{enumerate}
\item the union of the supports of $r(o)$ over any closed primal
simplex has diameter at most $2L$;
\item there is an augmentation-preserving equivariant chain map
\[
 \mathcal F:C_*(\mathcal O_L)\longrightarrow
 C_*^{\mathrm{dual}}(\widehat C)
\]
with finite output, such that if a dual cell to a primal simplex
$\tau$ occurs in $\mathcal F(e)$, then every label of every $r(o)$,
$o\in\tau$, is within $2L$ of every vertex of $e$.
\end{enumerate}
\end{lemma}

\begin{proof}
Use a sufficiently fine smooth triangulation, so that both primal
simplices and closed dual blocks have arbitrarily small mesh.  If all
images over a primal simplex share a label, every pair of labels in
their union is within $2L$.  This gives the first assertion.

We describe the chain construction explicitly.  The regular dual-cell
decomposition has closed cells that are balls and cellular subcomplexes.
Subdivide it barycentrically.  A simplex $s'$ of this subdivision is a
flag of dual cells; let $B(s')$ be the closure of its largest cell.
This is the unique minimal closed dual cell containing $s'$.  For a
face $t'$ of $s'$ we have $B(t')\subset B(s')$, so these closed balls
form a face-compatible acyclic carrier.

There is an augmentation-preserving chain map from the subdivision
chains to dual-cell chains carried by $B$.  To construct it, send a
vertex to a dual zero-cell in its carrier with coefficient one.  If
the map has been defined below dimension $j$, the image of the boundary
of a $j$-simplex is an augmented cellular cycle in its carrier, because
all face carriers are contained there.  Fill the cycle within this
contractible cellular subcomplex.  Make all choices on simplex orbits
and translate.  The deck action on the subdivision is free, so this is
equivariant.  Each carrier is finite, and hence each output is finite.

Now subdivide $\mathcal O_L$ equivariantly and approximate the actual
map $i$ closely by a simplicial map into the subdivision of the dual
decomposition. Choose the target stars inside sufficiently small evenly
covered balls. Simplicial approximation on the finite quotients can then
be joined to the original map by a homotopy staying in those stars.
Lifting that homotopy from $i$ gives an equivariant simplicial endpoint
that remains uniformly close to $i$ upstairs. Equivalently, one may choose
the stars directly on free vertex orbits. Compose source subdivision,
this simplicial chain map, and the carrier map just constructed.  This
defines $\mathcal F$.  Contractibility of $\mathcal O_L$ is not needed:
the only fillings take place inside the target carrier balls.

Consider a dual cell occurring on an original simplex $e$.  Its
carrier lies within a small mesh error of $i(z)$ for some $z\in e$.
The primal simplex $\tau$ dual to that cell meets it at its barycenter;
small primal mesh therefore puts every $o\in\tau$ within another small
error of such an $i(z)$.  Choose all meshes and the approximation close
enough to apply the common-label property to the pair $o,i(z)$.  A
common label $v$ belongs both to $r(o)$ and to $r(i(z))$.  Since
$ri=\pi$, the latter support consists of vertices of $e$.  For a label
$\gamma$ of $r(o)$ and a vertex $w$ of $e$, we obtain
\[
 d_\Gamma(\gamma,w)
 \le d_\Gamma(\gamma,v)+d_\Gamma(v,w)\le L+L.
\]
This proves the second assertion.
\end{proof}

\subsection{A compactly supported pairing}

For $0\le a\le L$ define
\[
 D_a=(\mathcal K_a\times\widehat C)/\Gamma,
 \qquad
 X=(\widetilde M\times\widehat C)/\Gamma.
\]
Both spaces project to $C$.  The oriented manifold $X$ has dimension
$n+d$, and the probability map induces the graph map
\[
 \begin{array}{ccc}
 X&\xrightarrow{\Phi_h}&D_L\\
 \downarrow&&\downarrow\\
 C&=&C
 \end{array}
 \qquad
 \Phi_h([\widetilde x,o])=[h(\widetilde x),o].
\]
This map is proper, as is the graph homotopy induced by the homotopy in
\cref{lem:probability-map}.  To see joint properness, cover compact $C$
by finitely many compact subsets contained in trivializing charts of
$\widehat C$.  In each chart a compact subset of $D_L$ has compact
first-factor projection.  Joint properness of the probability homotopy
controls the preimage in the $\widetilde M$ factor, and the parameter subset is compact.
A finite union gives the assertion.

The chain construction has a specific target: cap each simplex with the
detecting cocycle, retain the correct augmentation in degree $n$, and keep
the complementary dual cells close to the original labels. The following
lemma packages these properties before cellular duality is applied.

\begin{lemma}[A controlled cap operation]
\label{lem:controlled-cap}
There are equivariant linear maps with finite output
\begin{equation}
\label{eq:cap-chain-map}
 \begin{gathered}
 T_j:C_j(\mathcal K_L;\Q)\longrightarrow
 C_{j-n}^{\mathrm{dual}}(\widehat C;\Q),\qquad j\ge0,\\
 \partial T_j=(-1)^nT_{j-1}\partial,\qquad j\ge1,
 \end{gathered}
\end{equation}
where chain groups in negative degrees are zero. For every oriented
$n$-simplex $e$, the augmentation of $T_n(e)$ is $\omega(e)$. For every
oriented $j$-simplex $e$, if a dual cell to a primal simplex $\tau$ occurs
in $T_j(e)$, then every label of
$r(o)$, for every $o\in\tau$, is within $3L$ of every vertex of $e$.
\end{lemma}

\begin{proof}
On an oriented $j$-simplex $e=[v_0,\ldots,v_j]$ define the averaged
Alexander--Whitney diagonal by
\[
 \Delta_{\mathrm{AW}}e
 =\frac1{(j+1)!}\sum_{\varpi\in S_{j+1}}\sgn(\varpi)
   \sum_{k=0}^j
 [v_{\varpi(0)},\ldots,v_{\varpi(k)}]
 \otimes[v_{\varpi(k)},\ldots,v_{\varpi(j)}].
\]
Each ordering gives the usual diagonal on that simplex and its faces.
Restriction of a uniformly chosen ordering to a face is again uniformly
distributed, so the average commutes with the boundary. It is equivariant,
and both factors in every term are faces of $e$.

Apply $\omega$ to the first factor of degree $n$ and $\mathcal F\sigma$
to the second factor, and call the resulting graded map $T$. Its outputs
are finite. In the tensor boundary, the sum of terms from the first factor
vanishes after evaluation by $\omega$ because $\delta\omega=0$. The boundary in the
second factor carries the sign $(-1)^n$, proving
\eqref{eq:cap-chain-map}. When $j=n$, each contributing second factor is a
vertex. Its image under $\mathcal F\sigma$ has augmentation one, while
alternation makes every signed first-factor evaluation equal to $\omega(e)$.
Thus the augmentation of $T_n(e)$ is exactly $\omega(e)$.

Finally, every output of $T_j(e)$ comes from a face of $e$. By
\cref{lem:dual-carrier}, labels over its corresponding primal simplex are
within $2L$ of every vertex of that face. Since $e$ has diameter at most
$L$, they are within $3L$ of every vertex of $e$.
\end{proof}

Cellular duality will turn the finite output of $T$ in dual chains into primal
cochains. Its support bound will make the resulting product class compactly supported,
and its augmentation will recover the original cocycle pairing.

\begin{proposition}[Compact diagonal pairing]
\label{prop:compact-pairing}
There are a compact set $K_*\subset D_L$ and a class
\[
 u\in H^{n+d}(D_L,D_L\setminus K_*;\Q)
\]
such that
\begin{enumerate}
\item if $[t,o]\in K_*$, every label of $t$ is within $7L$ of
every label of $r(o)$;
\item the compactly supported evaluation
\begin{equation}
\label{eq:compact-pairing-nonzero}
 \ip{\Phi_h^*u,[X]_{\mathrm{lf}}}\ne0
\end{equation}
is nonzero.
\end{enumerate}
Here $[X]_{\mathrm{lf}}$ is the locally finite fundamental class.
\end{proposition}

\begin{proof}
Cellular duality on the oriented cover identifies finite dual chains
of degree $k$ with compactly supported primal cochains of degree
$d-k$.  Write this identification as $I$, choosing the degree signs
so that $\delta I=I\partial$.  At degree zero choose the sign for which
evaluation on the oriented locally finite fundamental cycle is the
augmentation; the remaining degree signs are then determined
successively by the boundary--coboundary identity.  On product cells,
define a cochain $U$ of degree $n+d$ by
\begin{equation}
\label{eq:product-cocycle}
 \begin{split}
 U(e^j,e'^{\,n+d-j})&=a_j I(T(e))(e'),\\
 a_n&=1,\qquad
 a_j=(-1)^{j+n+1}a_{j-1}\quad(j>n),
 \end{split}
\end{equation}
and set it to zero outside the indicated chain degrees.  The cochain
is invariant under simultaneous translation and thus descends to
$D_L$.  The product-cell structure is valid even if the action on
$\mathcal K_L$ has inversions: the freely lifted second-factor cells
give trivial stabilizers of product cells.

On a product cell of bidegree $(j,n+d+1-j)$, the two terms in
$\delta U$ are $a_{j-1}IT(\partial e)$ and
$(-1)^ja_jI\partial T(e)$.  By \cref{eq:cap-chain-map}, their
coefficients relative to $IT(\partial e)$ are
$a_{j-1}$ and $(-1)^{j+n}a_j$, which cancel.  Thus $U$ is a cocycle.

Call a product cell \emph{far} if all distances between vertices of
its first simplex and labels of $r$ over its closed second simplex
are greater than $4L$.  This condition is inherited by faces and is
translation invariant, so the far cells form a subcomplex
$D_{\mathrm{far}}\subset D_L$. The support bound in
\cref{lem:controlled-cap} shows that
$U$ vanishes on it.  Let $K_*$ be the union of the closures of all
non-far cells in the quotient.  There are only finitely many such
cells: after fixing a second-factor simplex from finitely many orbit
representatives, non-farness puts some first-factor vertex within
$4L$ of a label in a fixed finite set; the other vertices lie within
$L$ of that vertex.  Local finiteness of the word metric gives only
finitely many choices.

The same argument gives the required support bound.  One pair of
labels in a non-far product is within $4L$, the first simplex has
diameter at most $L$, and the union of second-factor labels has
diameter at most $2L$.  Every pair in its closure is consequently
within $4L+L+2L=7L$.  Furthermore
$D_L\setminus K_*\subset D_{\mathrm{far}}$, so the relative cellular
class of $U$ induces the required class $u$.

It remains to compute the pairing.  By proper homotopy invariance of
compactly supported cohomology, replace $h$ by $h_0$ and write
$\Phi_0$ for the resulting graph map.  On the product cell structure
of $X$, each open cell maps to an open product cell, possibly of
smaller dimension.  Thus
$X_{\mathrm{far}}=\Phi_0^{-1}(D_{\mathrm{far}})$ is a subcomplex.
Every cell not in it lies in the compact set $\Phi_0^{-1}(K_*)$, and
there are finitely many such cells by local finiteness.  The locally
finite fundamental class maps successively to
\[
 H_{n+d}(X,X\setminus\Phi_0^{-1}(K_*);\Q)
 \longrightarrow H_{n+d}(X,X_{\mathrm{far}};\Q).
\]
In the latter group it is represented by the finite product
fundamental cellular chain modulo $X_{\mathrm{far}}$.  Its coefficient
on each remaining top cell is the local orientation coefficient,
as follows by restricting further to local homology at an interior
point.  Hence evaluation by the relative cellular cocycle computes
the desired compact pairing.

Only bidegree $(n,d)$ contributes.  For each oriented fundamental
top simplex of $M$, choose one lift and sum over all oriented
fundamental top simplices of $\widehat C$.  The latter sum evaluates
$IT(e)$ on the locally finite parameter fundamental cycle and equals
the augmentation of $T(e)$, namely $\omega(e)$.  Collapsed labelled
top simplices contribute zero.  The total is exactly the nonzero
evaluation from \cref{lem:probability-map}, proving
\cref{eq:compact-pairing-nonzero}.
\end{proof}

\subsection{Distance profiles}

The compact pairing detects the original fundamental class. We next encode
its supporting region by finitely many scalar coordinates whose fibre
derivatives have a scale-independent bound. Their row count will determine
the dimension of the final parameter manifold only after these estimates
have been fixed.

Choose a finite cover of $C$ by small trivializing charts $U_\lambda$,
with $d+1$ colors and with charts of each color pairwise disjoint.
One obtains this by taking open stars of the vertices of a sufficiently
fine barycentric subdivision and coloring them by the dimensions of
the original simplices.  Shrink this finite cover and choose smooth
functions $0\le w_\lambda\le1$, supported in $U_\lambda$, so that at
least one equals one at each point of $C$.  On every sheet over a chart
choose an anchor $g_\lambda(o)\in\Gamma$, constant on the sheet and
equivariant under left translation, which is a common label of $r$
over that sheet.  Thus it is within $L$ of every label of every $r(o)$
there.

For each color and each $b\in B_\Gamma(1,1000L)$ use one row of a
vector $V_o(t)$.  Outside the charts of that color its value is zero;
on the sheet of a chart $U_\lambda$ it is
\begin{equation}
\label{eq:distance-profile-row}
 w_\lambda(\bar o)\sum_{\gamma\in\Gamma}t_\gamma
 \bigl(d_\Gamma(\gamma,g_\lambda(o)b)
             -d_\Gamma(g_\lambda(o),g_\lambda(o)b)\bigr),
 \qquad \bar o=[o].
\end{equation}
The disjointness within each color makes this definition unambiguous.
Compact support of each cutoff inside its chart makes it smooth in
the parameter wherever the probability coordinates are fixed.  Left
invariance of the word metric makes the vector invariant under
simultaneous translation, and hence it is defined on $D_L$.

Pad the vector by zero rows if necessary so that its number $q$ of
coordinates exceeds $n$.  Since the word balls have at most exponential
growth and $d\le2\abs{B_\Gamma(1,L)}+2$, there is a constant $J_1$,
independent of $R$, such that
\begin{equation}
\label{eq:profile-count}
 \log q\le J_1R.
\end{equation}
Choose an even integer $p_R\ge2$ with $q^{1/p_R}\le3$.  The norm
$\norm{\cdot}_{p_R}$ is smooth off zero and satisfies
\begin{equation}
\label{eq:profile-norms}
 \norm{z}_\infty\le\norm{z}_{p_R}\le3\norm{z}_\infty.
\end{equation}

\begin{lemma}[Profile estimates]
\label{lem:distance-profiles}
The map $V:D_L\to\R^q$ has the following properties.
\begin{enumerate}
\item It is continuous and proper, and
$\norm{V}_{p_R}\le24L$ on $K_*$.
\item If $t,t'\in\mathcal K_s$ satisfy
\[
 \max\{\norm{V_o(t)}_\infty,\norm{V_o(t')}_\infty\}\le110L,
 \qquad
 \norm{V_o(t)-V_o(t')}_\infty\le14R,
\]
then the union of their supports has diameter at most $L$.
\item The map
$\mathcal V=V\Phi_h:X\to\R^q$ is smooth.  With $K_0=sH$,
each of its rows satisfies
\[
 \sup_X\abs{d_{\widetilde M}\mathcal V_i}\le K_0,
\]
where $K_0$ is independent of $R$.
\end{enumerate}
\end{lemma}

\begin{proof}
Continuity follows from the coordinate definition and local finiteness.
At any point choose a chart whose cutoff equals one, and use its
$b=1$ row.  This row is the nonnegative average distance to its
anchor $g$.  For $t\in\mathcal K_L$ and any
$\gamma\in\supp t$, the triangle inequality gives
\[
 d_\Gamma(\gamma,g)
 \le\sum_{\gamma'}t_{\gamma'}d_\Gamma(\gamma',g)+L
 \le\norm{V_o(t)}_\infty+L.
\]
Over finitely many trivializing charts, a bound on $V$ therefore
restricts the first-factor labels to finitely many possibilities.
Closed inverse images of compact subsets are consequently compact.
On $K_*$, every input label is within $7L+L=8L$ of the active anchor.
Each distance difference in \cref{eq:distance-profile-row} is bounded
by that distance to the anchor, giving
$\norm{V}_\infty\le8L$, and hence the stated $24L$ bound.

For the second assertion, use a chart of weight one common to both
probabilities and let $g$ be its anchor.  Every label of either
probability is within $110L+s$ of $g$.  In particular a label
$\gamma\in\supp t$ equals $gb$ for one of the rows being used.
In this row the average of $d_\Gamma(\,\cdot\,,\gamma)$ over $t$
is at most $s$.  The row comparison implies that the average over
$t'$ is at most $s+14R$.  Since $\supp t'$ has diameter at most $s$,
every $\gamma'\in\supp t'$ satisfies
\[
 d_\Gamma(\gamma',\gamma)\le14R+2s<L.
\]
This applies to every label of $t$; the within-support distances
already satisfy the bound.

Finally, each row of $\mathcal V$ is locally a finite smooth sum.
In a $\widetilde M$ direction $\sum_\gamma dh_\gamma=0$.  Choose one
label from the current support and subtract its distance profile
before differentiating.  The remaining profile values have
oscillation at most $s$, because the distance function to any fixed
label is $1$-Lipschitz for $d_\Gamma$.  The derivative is thus bounded
by $s\abs{dh}_{\ell^1}\le sH$.  A smooth nonnegative coordinate that
vanishes at the point has zero differential, so zero probability
coordinates introduce no additional labels into this estimate.
\end{proof}

\subsection{From a relative pairing to local degree}

We now have the row count $q$ and its logarithmic bound. Relative duality
will turn the compact pairing into a transverse parameter cycle of
dimension $q-n$. Its intersections with the graph of the coordinate map
will supply the local degree, with different boundary pieces protecting
the two cycles.

Let $E=C\times\R^q$.  The set $V_o(\mathcal K_s)$ depends only on
$\bar o=[o]$, and we define
\[
 \rho(\bar o,y)=\dist_\infty\bigl(y,V_o(\mathcal K_s)\bigr).
\]
This is continuous.  In a trivialization, fixed witnesses give upper
semicontinuity.  For the reverse inequality, a sequence of almost
minimizing witnesses for bounded $y$ has bounded $V$-values; properness
in \cref{lem:distance-profiles} gives a convergent subsequence.  Its
limit is still in the closed subcomplex $\mathcal K_s$, and gives
lower semicontinuity.

Smooth $\rho$ within error $R/10$ on a neighborhood of
$\norm{y}_{p_R}\le110L$, obtaining $\rho_{\mathrm{sm}}$.  By Sard's
theorem \cite{Sard1942}, choose regular transverse levels
\begin{equation}
\label{eq:triad-levels}
 100L<T_0<101L,
 \qquad 5R<d_0<6R,
\end{equation}
also choosing $T_0$ regular for $\norm{\mathcal V}_{p_R}$ on $X$.
The latter function is smooth at these positive levels.  Then
\begin{equation}
\label{eq:compact-triad}
 P=\{(\bar o,y):\norm{y}_{p_R}\le T_0,
                         \ \rho_{\mathrm{sm}}(\bar o,y)\le d_0\}
\end{equation}
is a compact oriented manifold with corners.  Its radial and angular
sides are denoted by $E_R$ and $E_B$, respectively.  They satisfy
\begin{equation}
\label{eq:triad-gaps}
 \rho>4R\ \text{on }E_B,
 \qquad \rho<7R\ \text{throughout }P.
\end{equation}

There is a map $\psi:P\to D_L$ over $C$ with
\begin{equation}
\label{eq:witness-approximation}
 \norm{V\psi(\bar o,y)-y}_\infty<7R.
\end{equation}
Indeed, at each point choose a witness in $\mathcal K_s$ satisfying
this inequality, and continue it as a constant probability vector in
a small sheet trivialization.  It remains a witness on a neighborhood.
A finite partition of unity gives a convex mixture of these local
witnesses, interpreted in the common fiber.  Each witness has profile
norm at most $T_0+7R<110L$, and two witnesses at the same point have
profiles differing by less than $14R$.  By
\cref{lem:distance-profiles}, all labels in their union have pairwise
distance at most $L$.  Thus every mixture belongs to
$\mathcal K_L$.  Linearity of $V$ in probability coordinates proves
\cref{eq:witness-approximation}.

On $E_R$ its image satisfies
\[
 \norm{V\psi}_{p_R}\ge T_0-21R>24L,
\]
so $\psi(E_R)$ avoids $K_*$.  Consequently
\[
 u_P=\psi^*u\in H^{n+d}(P,E_R;\Q)
\]
is defined.  The compact domain
\[
 N=\{x\in X:\norm{\mathcal V(x)}_{p_R}\le T_0\}
\]
has smooth boundary and maps by
\[
 j=(\bar o,\mathcal V):(N,\partial N)\longrightarrow(P,E_R).
\]
Its image is disjoint from $E_B$, since $\rho$ is zero there.  At a
point of $N$, add $h(\widetilde x)$ to the family of witnesses used
by $\psi j$.  The same support test shows that linear interpolation
between $\psi j$ and $\Phi_h|_N$ stays in $D_L$.  On $\partial N$,
the profile norm throughout this homotopy is at least $T_0-21R>24L$.
Since $\Phi_h^{-1}(K_*)$ is in the interior of $N$, restriction and
excision, followed by this homotopy, give
\begin{equation}
\label{eq:relative-pairing}
 \ip{u_P,j_*[N,\partial N]}\ne0.
\end{equation}

\Cref{fig:degree-triad} records the complementary boundary conditions
used in the following duality argument.

\begin{lemma}[A transverse parameter cycle]
\label{lem:relative-degree}
There is a compact oriented smooth $(q-n)$-manifold $A_+$ with
boundary, together with a smooth map
\[
 (o_0,y_0):A_+\longrightarrow C\times\R^q,
\]
such that $\norm{y_0}_{p_R}<T_0$ everywhere,
$\rho(o_0,y_0)>4R$ on $\partial A_+$, and its intersections with
$j:N\to C\times\R^q$ are transverse, interior, and have nonzero
total oriented count.  The component maps may be taken constant in
the inward collar coordinate near $\partial A_+$.
\end{lemma}

\begin{proof}
Poincar\'e--Lefschetz duality for the triad
$(P;E_R,E_B)$ \cite[Theorem~3.43]{Hatcher2002} sends $u_P$ to a class
\[
 \alpha\in H_{q-n}(P,E_B;\Q),
\]
because $\dim P=d+q$.  Its relative intersection with
$j_*[N,\partial N]$ is, up to the fixed orientation convention,
\cref{eq:relative-pairing}.

We spell out the relative realization using the absolute theorem of
Thom \cite[Theorem~III.4]{Thom1954}.  Write $k=q-n$ and clear
denominators in $\alpha$.  Choose an integral relative cycle $z$ in
$(P,E_B)$ representing the resulting class.  Form the double
$D=P\cup_{E_B}P'$ along the angular face.  Compatible product collars
at $E_R\cap E_B$ give $D$ the structure of a smooth compact manifold
with boundary; the seam $E_B$ is a separating hypersurface meeting the
remaining boundary neatly.  If the seam is empty, this double is the
disjoint union of the two copies.  Write $z'$ for the reflected chain.
Because the boundary chains on the seam agree, $z-z'$ is an absolute
cycle in $D$.  The restriction and excision map
\[
 H_k(D;\Z)\longrightarrow H_k(D,P';\Z)
       \cong H_k(P,E_B;\Z)
\]
sends $[z-z']$ to $[z]$.

The compact manifold $D$ has finite polyhedral type. Transfer the doubled
class to a finite polyhedron homotopy equivalent to $D$, apply Thom's
theorem there, and compose back to $D$. This supplies a closed oriented
smooth $k$-manifold $B$ and a map
$f:B\to D$ whose fundamental class maps to a nonzero integer multiple
of $[z-z']$.  Reverse the orientation if necessary to make this
multiple positive.  Push $f$ away from the remaining boundary of $D$
by an inward collar deformation, chosen tangent to the seam.  Smooth
$f$ and perturb it transversely to the seam.  These operations are
homotopies in $D$, and compactness keeps the image away from the
boundary of the seam.  Thus $Q=f^{-1}(E_B)$ is a closed smooth
hypersurface, and $B_+=f^{-1}(P)$ and $B_-=f^{-1}(P')$ are compact
oriented smooth manifolds with common boundary $Q$.  Give
$B_+$ the orientation restricted from $B$.  The map
$H_k(B;\Z)\to H_k(B,B_-;\Z)\cong H_k(B_+,Q;\Z)$ sends $[B]$ to
$[B_+,Q]$: collar excision and the local orientation classes on the
interior of $B_+$ characterize this restriction.  Components wholly
in $B_-$ vanish in the relative group.  Naturality for the map of
pairs $(B,B_-)\to(D,P')$, followed by the displayed target
excision map, therefore gives
\[
 (f|_{B_+})_*[B_+,\partial B_+]
   =\text{a positive integer multiple of }\alpha
       \quad\text{in }H_k(P,E_B;\Q).
\]
This proves the required relative realization without prescribing
boundary data or requiring the representing map to be an embedding.

Write $a_+:B_+\to P$ for this representative and push its image off
$E_R$ by a collar deformation respecting $E_B$. At a corner one uses a
field pointing inward from the radial side and tangent to the angular
side; transversality of the two defining functions supplies such a field.
Compactness now separates $a_+(B_+)$ from $E_R$, and already separates
$j(N)$ from $E_B$. Consequently all coincidences of $a_+$ and $j$ lie in
a compact subset of the source interiors, with target values in
$\operatorname{int}P$.

Perturb $a_+$ only near this compact coincidence set. More precisely,
choose finitely many target coordinate neighborhoods with closures in
$\operatorname{int}P$ and use small flows of vector fields supported there
to make $a_+\times j$ transverse to the diagonal in $P\times P$.
Compactness permits the perturbation to be small enough that it creates
no coincidences outside these neighborhoods. The whole homotopy remains
a map of pairs $(B_+,\partial B_+)\to(P,E_B)$, is fixed near
$\partial B_+$, and keeps the image separated from $E_R$. Thus it
preserves the relative class and its pairing with $j_*[N,\partial N]$.
The dimension sum is $(q-n)+(n+d)=q+d$, so the transverse intersections
are isolated and have the asserted nonzero signed count. Set $A_+=B_+$
and write $(o_0,y_0)$ for the perturbed map into $E=C\times\R^q$.

Finally take a sufficiently small boundary collar, free of
intersections, and precompose the map there by a smooth collar
retraction that is constant in the inward collar coordinate near the
boundary and is the identity away from a slightly larger collar.
The image on that collar stays in a region where the strict
separation inequalities hold.  The resulting map is smooth, has the
same interior intersections, and has the claimed collar form.
\end{proof}

\begin{figure}[tb]
\centering
\begin{tikzpicture}[x=1cm,y=1cm,>=Latex,font=\small]
 \fill[blue!3] (-3,-1.55) rectangle (3,1.55);
 \draw[very thick,blue!65!black] (-3,-1.55)--(-3,1.55);
 \draw[very thick,blue!65!black] (3,-1.55)--(3,1.55);
 \draw[very thick,orange!75!black] (-3,-1.55)--(3,-1.55);
 \draw[very thick,orange!75!black] (-3,1.55)--(3,1.55);
 \draw[thick,->] (-3,-.35) .. controls (-1.8,.5) and (1.8,-.5) .. (3,.35);
 \draw[thick,->,purple!75!black] (-.45,-1.55) .. controls (.8,-.5) and (-.8,.5) .. (.45,1.55);
 \fill (0,0) circle (2pt);
 \node[above right] at (.15,.1) {intersection};
 \node[below] at (-1.9,-.03) {$j(N)$};
 \node[left,purple!75!black] at (-.45,-1) {$A_+$};
 \node[left,blue!65!black] at (-3,0) {$E_R$};
 \node[right,blue!65!black] at (3,0) {$E_R$};
 \node[above,orange!75!black] at (1.9,1.55) {$E_B$};
 \node[below,orange!75!black] at (1.9,-1.55) {$E_B$};
 \node at (-2.1,1.05) {$P$};
\end{tikzpicture}
\caption{A schematic transverse section of the compact triad.
The relative cycle $j(N)$ has boundary on the radial side $E_R$ and avoids
$E_B$. Its dual representative $A_+$ has boundary on $E_B$ and is pushed
off $E_R$. Their nonzero intersection becomes a local degree after the
parameter is pulled back and doubled. The drawing records incidence and
boundary conditions, not the dimensions of the cycles.}
\label{fig:degree-triad}
\end{figure}


\begin{proof}[Proof of \cref{thm:comparison-data}]
Double $A_+$ along its boundary, reversing the orientation on the
second copy, to obtain a closed oriented manifold $A$.  Reflect both
component maps $o_0,y_0$ across the seam.  The collar conclusion of
\cref{lem:relative-degree} makes the resulting maps smooth.  Pull back
$\widehat C\to C$ by $o_0:A\to C$, obtaining $\widehat A\to A$, and
form
\[
 Z=X\times_C A=(\widetilde M\times\widehat A)/\Gamma.
\]
The projection $X\to C$ is a submersion, so $Z$ is a smooth
$q$-manifold even though $o_0$ need not be a submersion.  Local sheet
trivializations identify it with a product of the
$\widetilde M$ fibre directions and the $A$ directions, with their product
orientation.  This also exhibits it as the stated cover of $M\times A$.
Continue to write $\mathcal V$ for the pulled-back profile map.

Choose once and for all a smooth function $\eta:[0,\infty)\to[0,1]$
equal to one on $[0,2]$ and zero on $[3,\infty)$, and define
\begin{equation}
\label{eq:comparison-map}
 G=\eta\bigl(\norm{\mathcal V}_{p_R}/T_0\bigr)
       \frac{\mathcal V-y_0}{R}.
\end{equation}
Properness of the profile map and compactness of $A$ imply that its
support is compact.  It is smooth, including at $\mathcal V=0$,
because the cutoff is constant in a neighborhood of that locus.
Write $a:Z\to A$ for the parameter projection, and let
\[
 \Omega=\{z\in Z:\ a(z)\in\operatorname{int}A_+,
                       \norm{\mathcal V(z)}_{p_R}<2T_0\}.
\]
This is relatively compact.  On its angular boundary the defining
boundary inequality for $A_+$ gives
$\norm{\mathcal V-y_0}_\infty>4R$.  On its other boundary,
\[
 \norm{\mathcal V-y_0}_{p_R}
 \ge2T_0-\norm{y_0}_{p_R}>T_0,
\]
and hence $\norm{\mathcal V-y_0}_\infty>T_0/3$.
The cutoff equals one on both parts of the boundary.  In view of
$T_0>100L=10000R$, these bounds prove the boundary assertion in
\cref{eq:comparison-degree}.

The zeros in $\Omega$ are exactly the intersections furnished by
\cref{lem:relative-degree}: if $\mathcal V=y_0$, their common
$p_R$-norm is less than $T_0$, so the point belongs to $N$.
We check orientations to identify their signed count.  In oriented
local lifts, the difference of the two map differentials, in tangent
variables $(v_M,v_C,v_A)$ and target rows $(C,\R^q)$, is
\[
 \begin{pmatrix}
  0&I&-do_0\\
  d_{\widetilde M}\mathcal V&d_C\mathcal V&-dy_0
 \end{pmatrix}.
\]
Substitute $v_C=v'_C+do_0(v_A)$.  This substitution has determinant
one.  Moving the $C$ variables to the front introduces only the fixed
dimensional sign $(-1)^{nd}$, and eliminating their identity block
leaves the differential of the pulled-back $\mathcal V-y_0$ on $Z$.
Thus the intersection signs agree with its zero signs up to one
fixed sign, independent of the point, sheet, and component.  Division
by the positive number $R$ does not change them.  Their nonzero sum
is therefore $\deg(G,\Omega,0)$ up to that sign.  Doubling the parameter
may cancel a global intersection count between its two copies; it
does not change the local degree over the first copy used here.

It remains to check the derivative bound.  By
\cref{lem:distance-profiles,eq:profile-norms},
\[
 \abs{d_{\widetilde M}\mathcal V_i}\le K_0,
 \qquad
 \abs{d_{\widetilde M}\norm{\mathcal V}_{p_R}}\le3K_0
 \quad(\mathcal V\ne0).
\]
In the cutoff transition region,
$\norm{\mathcal V-y_0}_\infty<4T_0$.  Writing
$E_\eta=\sup\abs{\eta'}$, differentiation of
\cref{eq:comparison-map} gives for every row
\[
 \abs{d_{\widetilde M}G_i}
 \le\frac{K_0}{R}
       +\frac{3K_0E_\eta}{T_0}\frac{4T_0}{R}
 \le\frac{K_0(1+12E_\eta)}{R}.
\]
The estimate also holds where the cutoff is constant.  Together with
\cref{eq:profile-count}, it proves \cref{eq:comparison-bounds} with
$J=\max\{J_1,K_0(1+12E_\eta),1\}$, increasing $R_0$ if necessary.
Finally $G$ is a smooth compactly supported map on the lifted product
geometry of the cover $Z\to M\times A$; each of its covariant
derivatives therefore has finite supremum for this fixed $R$.
\end{proof}
